\section{Introduction}
SQL deliberately hides physical execution choices.  That abstraction is right for applications, but it becomes dangerous when the hidden choice is the object being compared in a database study.  A benchmark author may ask whether two engines support pushdown, a system integrator may compare the evidence available in plan output, and a database researcher may group systems by join or adaptivity labels.  The same SQL query can then expose different public behavior: a local predicate rewrite, connector-side execution, a service-stage graph, a runtime filter, or only a logical explanation.  Before any latency or plan-quality number is interpreted, the comparison needs to know whether those public query-processing behaviors are being treated as the same object.

The problem is easiest to see in a concrete audited case.  For probe P0009, a keyword view equates four source records under a pushdown/explain label: BigQuery federated filter delegation, BigQuery query-stage evidence, Trino connector pushdown, and Trino runtime dynamic filtering.  The public contracts do not say the same thing: they distinguish connector-source delegation, service-stage evidence, aggregate/projection pushdown, and runtime filtering.  A paper that keeps only the keywords has not found an optimizer bug; it has lost the denominator needed to decide whether two public claims were comparable in the first place.

This is not a straw-man vocabulary.  We audit 26 public comparison surfaces drawn from engine documentation, plan interfaces, and optimizer-control descriptions.  All expose keyword-style or operator-family labels, twelve expose yes/no-style controls, and none exposes the full reference payload used by OptSem-C.  The work therefore asks a narrow systems question: when public optimizer behavior is summarized at the vocabulary level exposed by comparison surfaces, which source-supported distinctions disappear, and which retained fields make the comparison auditable again?

Database research already has precise abstractions for the objects it chooses to study.  Relational semantics starts from Codd's model~\cite{codd1970relational}; Selinger-style optimizers separate access-path search from costing~\cite{selinger1979access}; Volcano exposes rule-driven and parallel query processing~\cite{graefe1990volcano}; the Volcano optimizer-generator account makes extensible search explicit~\cite{graefe1993volcano}; and Cascades turns that line into a task-driven framework~\cite{graefe1995cascades}.  Chaudhuri surveys optimizer architecture~\cite{chaudhuri1998overview}; Ioannidis surveys cost and cardinality assumptions~\cite{ioannidis1996query}; and Steinbrunn et al. study join-order search heuristics~\cite{steinbrunn1997heuristic}.  Public optimizer behavior is usually compared with much coarser vocabularies.  A cell labeled ``pushdown'' may mean predicate movement inside the engine, partition pruning, aggregation delegation, join delegation, limit delegation, or runtime filtering.  A cell labeled ``join'' may refer to a logical operator, an estimated physical operator, a distributed exchange, or a runtime profile entry.  A claim that two engines are ``cost based'' may hide whether statistics are local, connector-provided, unavailable, or revised during execution.  These shortcuts are not harmless summarization: they can create projection-induced contract collisions (a coarse vocabulary declares equality among public contracts whose reference signatures differ).

OptSem-C treats this precision problem as a query-processing systems problem, not as prose cleanup.  A \emph{public optimizer contract} is the optimizer behavior an engine exposes through plan interfaces, tuning surfaces, and system descriptions.  A source record can support a contract that is not visible in a particular textual plan, and a plan tag can expose behavior without proving the private cost model.  The object is therefore not a hidden implementation trace and not a latency measurement.  OptSem-C represents each public statement as a guarded contract rule over query features.  The rule records the action, the guard under which it applies, where the action is placed, which optimizer layer owns it, when the decision is made, how the behavior is observable, and whether the public contract requires, permits, forbids, or does not specify the action.

\textbf{Research questions.}  We ask four finite, auditable questions.  RQ1 asks whether public optimizer claims can be encoded as source-linked contracts and executable probes without becoming a latency benchmark.  RQ2 asks how often comparison vocabularies that appear in public feature claims collapse distinct source-supported contracts.  RQ3 asks which semantic fields eliminate those collisions, and which robustness checks delimit the repair claim.  RQ4 asks whether the projection-audit inner loop and SQL-denominator checks are cheap enough to rerun when the public source denominator changes, rather than trusted as a one-off manual inspection.

The central result is a projection-kernel account of these contract collisions.  Comparing optimizer contracts through a keyword or operator vocabulary applies a predeclared projection to reference signatures: the full public evidence payload inside the declared schema, not an oracle over private optimizer state.  If reference signatures disagree but fall into the same projected class, the projection has declared an equivalence that source-supported public signatures do not support.  Contract collisions are therefore denominator errors for the comparison itself; performance or cost measurements should be interpreted only after the compared behavior has been named.

OptSem-C turns that denominator question into a finite comparison relation with executable counter-witnesses, exact repair certificates, and negative controls.  The contract schema, source manifest, projection portfolio, motif crosswalk, and repair-field universe are frozen before witness counting; execution checks validate probes but never admit rules or choose repairs.  Layer and placement are therefore not learned explanations for the failures: they are schema dimensions inherited from optimizer architecture and federation literature, then tested in the same predeclared field universe as every other retained field.  Keyword projection declares 2,284 engine--probe equivalences: 2,030 have matching reference signatures and 254 are separated by reference signatures.  Operator-only projection declares 2,268 equivalences, including 238 collisions; the yes/no projection contributes six sparse source-supported records.  Across the three vocabularies, the total is 498 projection-witness records over four observed engine pairs and 486 distinct probes, where a record is keyed by projection, engine pair, and probe.  Restoring placement repairs keyword collisions; restoring optimizer layer repairs operator-only collisions; retaining layer and placement is one interpretable two-field repair basis for all headline records within the admitted public-contract corpus.

The robustness checks are designed to block three tempting overclaims.  First, the probe set is a rule-aware coverage benchmark, not an independent population sample, so we do not report statistical generalization beyond the frozen denominator.  Second, source-removal is an information-deletion stress rather than a new sample: keyword and operator-only remain nonzero after every single-source removal, and the audit separately records retained original witnesses and newly induced witnesses.  The sparse yes/no projection is reported as source-sensitive and stays outside the broad robustness claim.  Third, repair fields are not learned from failed engine-pair transfer: the reported layer-and-placement basis resolves feature-family and engine-family stress splits, while point-learned engine-pair repairs fail on held-out pairs and remain a boundary.  A field-only stress test over all 256 retained-field vocabularies leaves only 44 unsafe regions, all covered by concrete counter-witnesses; after excluding action-variant identifiers, no singleton semantic field is safe.

We also introduce OptSemBench-C, a covering benchmark for optimizer-contract behavior.  Its denominator is not raw query count.  It is coverage of optimizer-relevant feature interactions: source boundary, connector capability, join shape, predicate class, aggregation, statistics need, reuse, distribution, adaptivity, and control surface.  The generated probes are executable SQL representatives.  The probe space is then cross-checked against DPccp-style join enumeration~\cite{moerkotte2006dpccp}, JOB optimizer evidence~\cite{leis2015good}, later JOB diagnostics~\cite{leis2018looking}, exact-cardinality testing~\cite{chaudhuri2009exact}, progressive optimization~\cite{markl2004progressive}, mid-query reoptimization~\cite{kabra1998reoptimization}, adaptive query processing~\cite{deshpande2007adaptive}, and eddies~\cite{avnur2000eddies}.  We use this audit as a published-motif vocabulary stress test: 90 optimizer-motif requirements are covered by 71 distinct generated SQL representatives, so the check is concrete without pretending that generated representatives are copied JOB, TPC-H, or TPC-DS workloads.

\textbf{Contributions.} C1 defines a public optimizer contract as a query-processing object separate from SQL result equivalence and runtime performance; Section~\ref{sec:model} gives the guarded rule model and the state-join semantics.  C2 turns comparison validity into a finite projection audit: reference signatures define the public evidence relation, projected signatures expose counter-witnesses, and repair certificates identify fields that make a vocabulary safe.  C3 contributes OptSemBench-C as a denominator design for optimizer-contract comparisons across seven analytical SQL engines; Section~\ref{sec:benchcorpus} reports source-linked admission, executable probes, external-motif coverage, and zero-failure admission checks.  C4 reports the projection kernels, anti-overfitting checks, resource profile, and repair certificates: 498 headline projection-witness records, layer and placement as an interpretable repair basis for those records, and cloud replay evidence for the comparison inner loop.

Each central claim is tied to a reproducible record rather than to prose alone.  The model section defines the state table and projection obligations; the corpus section exposes source identifiers, line spans, and digest prefixes; the evaluation tables connect SQL execution, motif coverage, collision counts, and repair certificates to regenerated outputs.  The case studies then show how those aggregate certificates change comparison judgments.
